\documentclass{birkau}
\usepackage{amsmath,amssymb,url}
\usepackage{hyperref}

\newtheorem{theorem}{Theorem}[section]
\newtheorem{proposition}[theorem]{Proposition}
\newtheorem{lemma}[theorem]{Lemma}
\newtheorem{corollary}[theorem]{Corollary}
\theoremstyle{definition}
\newtheorem{definition}[theorem]{Definition}
\newtheorem{example}[theorem]{Example}
\theoremstyle{remark}
\newtheorem{remark}[theorem]{Remark}

\begin{document}

\title[Finite-complement congruence topology]{A Finite-Complement Congruence Topology for Partial Algebras}
\author{Adrian Puha}
\address{Independent researcher}
\subjclass[2020]{Primary 08A55; Secondary 08A30, 08B20}
\keywords{partial algebra, free completion, Rees congruence, congruence topology, profinite topology}

\begin{abstract}
Let $D$ be a partial algebra with carrier $A$, and let $F(D)$ be its free
compatible total extension.  We consider compatible total extensions that fix
$A$ pointwise and contain only finitely many elements outside $A$.  Their
kernels define a descending family of congruences on $F(D)$ and hence a
Hausdorff congruence topology.

For every finite subset of $F(D)$ there is one such extension on which the
canonical homomorphism is injective.  This is an elementary Rees-congruence
construction: retain its finite subterm closure and collapse the complement to
one class.  When $A$ is finite, the resulting topology is exactly the ordinary
profinite topology.

The case of interest is therefore an infinite retained base, where the allowed
target algebras may themselves be infinite.  We prove a sufficient condition
for the Hausdorff completion to be strictly larger than $F(D)$.  If $D$ has an
undefined base application and admits a nontrivial unary context whose orbits
on $A$ have uniformly bounded finite cardinality, then iteration from the
undefined term yields a Cauchy sequence with no limit in $F(D)$.  The identity
on the base is the simplest special case.  We also exhibit an infinite partial
algebra for which the induced topology is discrete, so the hypothesis is not
implied by partiality alone.
\end{abstract}

\maketitle

\section{Introduction}\label{sec:introduction}

Let $D=(A,\Sigma,\delta)$ be a one-sorted partial algebra with finite binary
signature.  We write $F(D)$ for the standard free compatible total extension of
$D$: applications already defined on the base are evaluated, while undefined
applications remain as normal terms.  Free completion of partial algebras is
classical and will be used only as background
\cite{ClimentVidalCosmeLlopez2023,Nesin2011}.
Related total-algebraic approaches encode partial algebras inside suitable total
algebraic structures rather than freely adjoining normal terms.  In particular,
Diaconescu gives a Horn encoding based on existence equality, with a satisfaction
condition and persistent-liberality properties \cite{Diaconescu2009}.  That
viewpoint is conceptually related to, but distinct from, the finite-complement
compatible extensions considered here.

We study compatible pairs $(T,i)$ in which $i:A\to T$ is injective and only
finitely many elements of $T$ lie outside $i(A)$.  We identify $A$ with its
image $i(A)$ throughout, and therefore write
\[
 A\subseteq T,
 \qquad |T\setminus A|<\infty.
\]
Every operation value already defined in $D$ has the same value in $T$ under
this identification.  The restriction measures how much genuinely new
structure a total extension must introduce: the base is kept intact, so the
finite complement counts only elements added beyond it.  The universal property
of $F(D)$ gives a canonical homomorphism
\[
 \pi_T:F(D)\longrightarrow T
\]
fixing $A$ pointwise.  We call such a pair a \emph{finite-complement compatible
extension}.  A one-point completion is the simplest example of this kind
\cite{Nesin2011}; here we use all finite-complement compatible extensions rather
than a distinguished sink completion.

The first observation is elementary.  If $X\subseteq F(D)$ is finite, then one
can retain the finite subterm closure of $X$ and collapse everything else into
one class.  The complement is absorbing under every basic operation, so its
collapse gives a Rees congruence in the usual universal-algebraic sense
\cite{Tichy1981}.  The resulting quotient is a finite-complement compatible
extension and the quotient map is injective on $X$.  This is the algebraic
version of the standard finite-subterm plus sink construction used in finite
tree automata \cite{ComonEtAl2008}.  We record it because it supplies the
separating congruences used below.

For $m\ge1$, let $\theta_m$ be the intersection of the kernels arising from
finite-complement compatible extensions with at most $m$ new elements.  The
relevant kernels form a set of congruences on $F(D)$, so this intersection is
set-indexed.  The separation observation implies
\[
 \bigcap_{m\ge1}\theta_m=\Delta_{F(D)}.
\]
Thus $(\theta_m)$ is a countable base of congruence entourages for a separated
uniformity on $F(D)$, in the standard sense of congruence topologies on
universal algebras \cite{BulmanFleming1971}.  Because the basis is countable and
nested, the uniformity is metrizable; consequently Cauchy sequences suffice to
describe its completion.

The finite-base case is not a new completion theory.  When $A$ is finite, the
finite-complement quotients that fix $A$ are cofinal among all finite quotients
of $F(D)$.  Consequently the topology is exactly the ordinary profinite
topology \cite{DaveyGouveiaHaviarPriestley2011,ChenAdamekMiliusUrbat2016}.
The genuinely different situation is therefore an infinite retained base,
where the allowed target algebras may themselves be infinite and the
congruences $\theta_m$ need not have finite index.

Our main result gives a sufficient condition for the completion to be strictly
larger than $F(D)$ in this setting.  Suppose $D$ has an undefined base
application, producing a non-base normal term $c_0$, and suppose there is a
nontrivial one-hole unary context $C[-]$ whose induced map $C^D:A\to A$ has
uniformly finite orbits: there is $B$ such that
\[
 \bigl|\{(C^D)^n(a):n\ge0\}\bigr|\le B
 \qquad(a\in A).
\]
The pointwise identity case is the simplest example, but the hypothesis also
covers constant, idempotent, and uniformly bounded periodic behaviour on the
base.

Iterate the context in the free completion, $c_{n+1}=C[c_n]$.  In every
extension with at most $m$ new elements, the interpreted orbit either remains
outside the base, where it has at most $m$ values, or enters the base and then
follows one of the uniformly bounded base orbits.  Hence the whole interpreted
orbit has at most $m+B$ values.  Factorially indexed iterates therefore
synchronize uniformly over all such extensions and form a Cauchy sequence for
every $\theta_m$.

The sequence does not converge to a term of $F(D)$.  For any proposed term
$x$, a Rees quotient retaining the subterms of $x$ and a sufficiently late
iterate sends the next iterate, and hence the whole subsequent tail, into the
collapsed class while leaving $x$ singleton.  The Hausdorff completion
therefore contains a new point.

The sufficient hypothesis is not merely a disguised form of partiality.  We
also give an infinite partial algebra for which $\theta_1$ is equality, so the
induced topology is discrete and already complete.  Thus genuine partiality
alone does not force new completion points over an infinite retained base.

The resulting uniformity should also be distinguished from the classical
prefix-depth ultrametrics on trees and their associated infinite-tree spaces
\cite{ArnoldNivat1980,Courcelle1983}.  We give an explicit comb sequence that
is Cauchy for prefix-depth agreement but is not Cauchy for the
finite-complement congruence uniformity.

The finite-complement condition is formally analogous to finite Rees index for
subsemigroups \cite{RuskucThomas1998,CainMaltcev2013} and to separability
conditions studied in \cite{MillerOReillyQuickRuskuc2022}.  The present setting
is different because the distinguished copy of $A$ need not be a subalgebra of
$F(D)$.  In fact, it is closed under the total operations of $F(D)$ if and only
if $D$ was already total: any undefined basic application produces a non-base
normal term immediately outside $A$.

The paper is organized accordingly.  Section~\ref{sec:free-completion} recalls
the free compatible extension.  Section~\ref{sec:finite-separation} gives the
Rees-congruence separation lemma.  Section~\ref{sec:completion} defines the
congruence topology, identifies the finite-base case with the profinite
topology, proves the infinite-base criterion, gives an arithmetic instance,
gives the discrete counterexample, and compares the uniformity with
prefix-depth agreement on trees.

The manuscript is self-contained.  Supplementary machine-checked developments
cover several underlying constructions, special cases, and the prefix-depth
contrast.

\section{Free compatible completion}\label{sec:free-completion}

We use the standard conventions for partial algebras; see, for example,
\cite{Burmeister1986}.  Let $D=(A,\Sigma,\delta)$ be a one-sorted partial
algebra with binary signature $\Sigma$.  For $f\in\Sigma$ and $a,b\in A$, the
value $\delta_f(a,b)$ is either an element of $A$ or is undefined.

Throughout, evaluation of composite terms in $D$ is \emph{strict}: a term
$f(s,t)$ is defined only when both $s$ and $t$ are defined and the corresponding
base application $\delta_f(s^D,t^D)$ is defined.  Thus, whenever a composite
term is defined in $D$, every subterm occurring in its evaluation is defined as
well.

A \emph{normal term} over $D$ is defined recursively.  Every $a\in A$ is a
normal term.  If $s,t$ are normal terms and $f\in\Sigma$, then $f(s,t)$ is a
normal term unless both $s,t$ lie in $A$ and $\delta_f(s,t)$ is defined; in
that case the application is replaced by the corresponding base value.  Let
$F(D)$ denote the set of normal terms.  Each basic operation becomes total on
$F(D)$ by forming the formal application and applying this normalization rule.

The inclusion $A\hookrightarrow F(D)$ is injective and preserves every value
already defined in $D$.  A total $\Sigma$-algebra $T$ with a map $i:A\to T$ is
\emph{compatible with $D$} when
\[
 \delta_f(a,b)=c
 \quad\Longrightarrow\quad
 f_T(i(a),i(b))=i(c).
\]

\begin{theorem}[Universal property]\label{thm:free-completion}
For every compatible pair $(T,i)$ there is a unique homomorphism
\[
 \pi_T:F(D)\longrightarrow T
\]
satisfying $\pi_T(a)=i(a)$ for every $a\in A$.
\end{theorem}

\begin{proof}
Evaluate normal terms recursively in $T$.  A non-base normal term has the form
$f(s,t)$, and its value is defined to be
$f_T(\pi_T(s),\pi_T(t))$.  If an application of $f$ to two base elements was
already defined in $D$, it is not a non-base normal term: it has already been
reduced to its base value.  Compatibility therefore makes recursive evaluation
respect the normalization rule.  The homomorphism property and uniqueness
follow by induction on term formation.
\end{proof}

\section{Finite-complement separation}\label{sec:finite-separation}

For a normal term $x\in F(D)$, let $\operatorname{Sub}(x)$ denote its finite
set of normal subterms.  If $X\subseteq F(D)$ is finite, set
\[
 S_X=\bigcup_{x\in X}\operatorname{Sub}(x),
 \qquad
 N_X=S_X\setminus A.
\]
Thus $N_X$ is finite and $S_X$ is closed under taking subterms.

\begin{lemma}[Finite-complement Rees quotient]
\label{lem:finite-complement-rees}
For every finite $X\subseteq F(D)$ there is a congruence $\rho_X$ on $F(D)$
such that every element of $A\cup N_X$ is a singleton $\rho_X$-class and, if
nonempty, $F(D)\setminus(A\cup N_X)$ is one additional class.  Consequently
$F(D)/\rho_X$ is a finite-complement compatible extension of $D$ and the
quotient map is injective on $X$.
\end{lemma}

\begin{proof}
Put
\[
 I_X=F(D)\setminus(A\cup N_X).
\]
We first observe that $I_X$ is absorbing in every coordinate: if one argument
of a basic operation lies in $I_X$, then so does the result.  Indeed, let
$f(s,t)$ be the result of applying a basic operation and suppose, for example,
that $s\in I_X$.  The result cannot lie in $A$, because reduction to a base
value is possible only when both arguments lie in $A$.  Nor can it lie in
$N_X$: if the normal term $f(s,t)$ were a selected subterm, then its subterm
$s$ would belong to $S_X$, contrary to $s\in I_X$.  The same argument applies
when $t\in I_X$.

Define
\[
 \rho_X=\Delta_{F(D)}\cup(I_X\times I_X).
\]
The absorption property makes $\rho_X$ a congruence; when $I_X$ is nonempty,
this is the Rees congruence determined by the absorbing subalgebra $I_X$
\cite{Tichy1981}.  Every base element remains a singleton class.  Outside the
image of $A$ the quotient has the singleton classes indexed by $N_X$, together
with at most one additional class represented by $I_X$.  Hence
\[
 |(F(D)/\rho_X)\setminus A|\le |N_X|+1<\infty.
\]
Since $X\subseteq A\cup N_X$, the quotient map is injective on $X$.
\end{proof}

The construction is the algebraic form of the familiar finite-subterm
construction with one sink state in tree automata \cite{ComonEtAl2008}.

\begin{corollary}[Separation]
\label{cor:finite-complement-separation}
If $x\ne y$ in $F(D)$, then there is a finite-complement compatible extension
$T$ for which the canonical map $F(D)\to T$ separates $x$ and $y$.
\end{corollary}

\begin{remark}
When $A$ is infinite, finite-complement extensions are not in general closed
under ordinary products over the diagonal copy of $A$: already the diagonal
copy of $A$ in $(A\sqcup\{\ast\})^2$ has infinitely many points outside it.
This observation is not needed in the proof above; it merely explains why the
usual product-of-separators argument does not stay inside the same class.
\end{remark}

\section{The induced congruence topology}\label{sec:completion}

For $m\ge1$, let $\mathcal K_m$ be the set of congruences on $F(D)$ that occur
as kernels of canonical maps
\[
 \pi_T:F(D)\longrightarrow T
\]
where $T$ is a finite-complement compatible extension with
$|T\setminus A|\le m$, and put
\[
 \theta_m=\bigcap\mathcal K_m.
\]
This is a set-indexed intersection: $\mathcal K_m$ is a subset of the set of
all congruences on $F(D)$.  Each $\theta_m$ is a congruence and
$\theta_{m+1}\subseteq\theta_m$.  By
Corollary~\ref{cor:finite-complement-separation},
\[
 \bigcap_{m\ge1}\theta_m=\Delta_{F(D)}.
\]
Hence $(\theta_m)_{m\ge1}$ is a countable base of congruence entourages for a
separated uniformity on $F(D)$, a standard construction for universal algebras
\cite{BulmanFleming1971}.  We denote its Hausdorff completion by
$\widehat{F(D)}_{\mathrm{fc}}$.  Since the uniformity is separated, the
canonical map
\[
 F(D)\longrightarrow\widehat{F(D)}_{\mathrm{fc}}
\]
is injective; the question below is whether it is surjective.

Because the congruence basis is countable and nested, this uniformity is
metrizable.  Explicitly, put $\theta_0=F(D)^2$ and, for $x\ne y$, let
\[
 r(x,y)=\max\{m\ge0:(x,y)\in\theta_m\},
 \qquad d(x,y)=2^{-r(x,y)},
\]
with $d(x,x)=0$.  Then $d$ is an ultrametric inducing the same uniformity.
Consequently it is enough to work with Cauchy sequences rather than general
Cauchy filters.

\subsection{Finite bases}

When $A$ is finite, the preceding topology is not a new kind of completion.

\begin{proposition}[Finite base gives the profinite topology]
\label{prop:finite-base-profinite}
Assume that $A$ and $\Sigma$ are finite.  Then the topology generated by the
congruences $\theta_m$ is the ordinary profinite topology on $F(D)$.
\end{proposition}

\begin{proof}
First, each $\theta_m$ has finite index.  Up to relabelling of the finitely many
points outside the fixed finite set $A$, there are only finitely many total
$\Sigma$-algebra structures on carriers of size at most $|A|+m$.  Thus only
finitely many kernels occur in the defining intersection for $\theta_m$, and
an intersection of finitely many finite-index congruences again has finite
index.

Conversely, let $\alpha$ be any finite-index congruence on $F(D)$.  Apply
Lemma~\ref{lem:finite-complement-rees} to the finite set $X=A$.  This gives a
finite-index congruence $\beta$ whose restriction to $A$ is equality.  Put
$\gamma=\alpha\cap\beta$.  Then $\gamma$ has finite index,
$\gamma\subseteq\alpha$, and $A$ remains pointwise distinct in
$F(D)/\gamma$.  Identifying the image of $A$ with $A$ itself, the quotient is a
finite-complement compatible extension.  If its complement has size $m$, put
$M=\max\{m,1\}$.  Then
\[
 \theta_M\subseteq\gamma\subseteq\alpha.
\]
Thus $(\theta_m)_{m\ge1}$ is cofinal among the finite-index congruences and
induces the ordinary profinite topology.
\end{proof}

For infinite $A$, an allowed target may itself be infinite and the congruences
$\theta_m$ need not have finite index.  The rest of the section concerns this
case.

\subsection{Uniformly finite base orbits}

We use one elementary observation about finite deterministic orbits.

\begin{lemma}[Factorial synchronization]\label{lem:finite-orbit-factorials}
Let $s_0,s_1,\ldots$ be an orbit of a self-map.  If the orbit takes at most
$N$ distinct values, then $s_{k!}$ is independent of $k$ for all sufficiently
large $k$.  The threshold can be chosen using only $N$.
\end{lemma}

\begin{proof}
Among the first $N+1$ positions two coincide.  The orbit is therefore eventually
periodic, with preperiod and period bounded by $N$.  Every sufficiently large
factorial lies beyond the preperiod and is divisible by the period.
\end{proof}

A \emph{one-hole unary context} $C[-]$ is a term with exactly one occurrence of
a distinguished variable, representing the hole, and no other variables.  We
call it \emph{nontrivial} if it is not the bare hole $[-]$.  Suppose that
$C[a]$ is defined in $D$ for every $a\in A$; by the strict evaluation
convention of Section~\ref{sec:free-completion}, it then induces a map
$C^D:A\to A$.

\begin{lemma}[Compatibility through contexts]\label{lem:context-compatibility}
Let $(T,i)$ be compatible with $D$.  If $C[a]$ is defined in $D$, then
evaluating $C[a]$ in $T$ gives $i(C^D(a))$.
\end{lemma}

\begin{proof}
Induct on the construction of the context.  Since definedness in $D$ is strict,
every operation node encountered in the evaluation has defined base arguments.
Compatibility of $T$ with $D$ therefore preserves the value at each node.
\end{proof}

We say that $C[-]$ has \emph{uniformly finite base orbits} if $C^D$ is defined
on all of $A$ and there is an integer $B\ge1$ such that
\[
 \bigl|\{(C^D)^n(a):n\ge0\}\bigr|\le B
 \qquad\text{for every }a\in A.
\]
The identity case has $B=1$; maps of uniformly bounded finite order,
idempotent maps, and constant maps are further elementary examples.

\begin{theorem}[Infinite-base completion criterion]
\label{thm:infinite-base-criterion}
Assume that $A$ is infinite.  Suppose that some basic application in $D$ is
undefined, and let $c_0\in F(D)\setminus A$ be the corresponding normal term.
If there is a nontrivial one-hole unary context $C[-]$ with uniformly finite
base orbits, then the canonical map
\[
 F(D)\longrightarrow\widehat{F(D)}_{\mathrm{fc}}
\]
is not surjective.
\end{theorem}

\begin{proof}
Define $c_{n+1}=C[c_n]$.  Since $C[-]$ is nontrivial and $c_0$ is non-base,
the occurrence of $c_n$ cannot be removed by the normalization rule: every
operation node on the path from the hole has a non-base argument.  Thus the
construction depth strictly increases and the $c_n$ are pairwise distinct.

Let $B$ bound the base orbits of $C^D$, and put $d_k=c_{k!}$.  Fix $m$ and a
finite-complement compatible extension $T$ with $|T\setminus A|\le m$.  The
images of the $c_n$ form an orbit under the unary operation induced by $C[-]$
on $T$.  If this orbit never enters $A$, it takes at most $m$ values.  If it
enters $A$, the part before first entry has at most $m$ distinct values outside
$A$; after entry, Lemma~\ref{lem:context-compatibility} identifies the remainder
with an orbit of $C^D$, which has at most $B$ values.  Hence in every case the
whole orbit has at most $m+B$ distinct values.

By Lemma~\ref{lem:finite-orbit-factorials}, for all sufficiently large $k,l$,
with a threshold depending only on $m+B$ and not on $T$,
\[
 \pi_T(d_k)=\pi_T(d_l).
\]
Therefore $d_k\equiv d_l\pmod{\theta_m}$ for all sufficiently large $k,l$.
Since $m$ was arbitrary, $(d_k)$ is Cauchy.

It remains to show that this Cauchy sequence has no limit in $F(D)$.  Fix
$x\in F(D)$.  Choose $N$ so large that the construction depth of $c_N$ exceeds
that of $x$, and apply Lemma~\ref{lem:finite-complement-rees} to
$X=\{x,c_N\}$.  The term $c_{N+1}=C[c_N]$ has greater depth than both selected
roots and hence is not among their selected subterms, so it lies in the
collapsed class.  Because that class is absorbing under every basic operation,
substitution into the context keeps all later $c_j$ in the same class.  The
term $x$, however, remains a singleton class.  Thus this one
finite-complement quotient separates $x$ from every sufficiently late $d_k$.
If the quotient has complement size $q$, put $M=\max\{q,1\}$.  Its kernel
belongs to $\mathcal K_M$, so for every sufficiently late $k$ we have
$(d_k,x)\notin\theta_M$.  Hence $(d_k)$ does not converge to $x$.

Since $x$ was arbitrary, $(d_k)$ has no limit in $F(D)$.  Its limit in the
Hausdorff completion is therefore outside the image of $F(D)$.
\end{proof}

\begin{corollary}[Pointwise-fixed base]\label{cor:base-fixing}
The conclusion of Theorem~\ref{thm:infinite-base-criterion} holds in particular
when $C^D(a)=a$ for every $a\in A$.
\end{corollary}

\begin{example}[Natural numbers]\label{ex:natural-numbers}
Let $D$ have carrier $\mathbb N$, with the usual addition and multiplication
and with a binary division operation that is undefined whenever the denominator
is zero; its values at nonzero denominators are immaterial here.  In $F(D)$ the
expression $0/0$ is therefore a non-base normal term.  The one-hole context
\[
 C[-]=[-]+0
\]
is nontrivial and satisfies $C^D(n)=n$ for every $n\in\mathbb N$.  Hence
Theorem~\ref{thm:infinite-base-criterion} applies: the completion of $F(D)$ for
the finite-complement congruence topology is strictly larger than $F(D)$.
\end{example}

The example uses partial division only to provide an undefined base
application; no numerical value is assigned to $0/0$.  This should be
distinguished from totalized approaches such as meadows and wheels
\cite{BergstraHirshfeldTucker2009,Carlstrom2004}.

The criterion is only sufficient.  The following example shows that this
qualification matters: genuine partiality alone does not force a nontrivial
completion when the retained base is infinite.  In this example no nontrivial
context is defined on the whole base, so the hypothesis of
Theorem~\ref{thm:infinite-base-criterion} fails in the strongest possible way.

\begin{proposition}[A discrete infinite-base example]
\label{prop:discrete-example}
Let $A$ be infinite, let the signature consist of one binary symbol $f$, and
let $f$ be nowhere defined on $A$.  Then $D$ is genuinely partial but
$\theta_1=\Delta_{F(D)}$.  Hence the finite-complement congruence topology on
$F(D)$ is discrete and its Hausdorff completion adds no new points.
\end{proposition}

\begin{proof}
Fix distinct $x,y\in F(D)$ and let $S$ be the finite subterm closure of
$\{x,y\}$.  Since $A$ is infinite, choose an injective map
$\lambda:S\to A$ that fixes every element of $S\cap A$.  Fix $a_0\in A$ and
define a total binary operation on $A$ by
\[
 f_A(\lambda(s),\lambda(t))=\lambda(f(s,t))
\]
whenever $f(s,t)\in S$, and put $f_A(r,s)=a_0$ on every remaining pair.
These prescriptions are consistent because $\lambda$ is injective.  As the
partial operation of $D$ is nowhere defined, there are no compatibility
constraints to violate.  Induction on subterms gives that the canonical map
$F(D)\to A$ agrees with $\lambda$ on $S$, and hence separates $x$ and $y$.
Thus every distinct pair is separated by an extension with zero-point
complement.  It follows that $\theta_1$ is equality, so the induced uniformity
is discrete and already complete.
\end{proof}

\begin{remark}[Contrast with prefix-depth completion]
\label{rem:prefix-depth-contrast}
The finite-complement uniformity is not the usual prefix-depth uniformity on
terms.  Suppose $\delta_f(a,z)$ is undefined and define a left comb by
$c_0=a$ and $c_{k+1}=f(c_k,z)$.  For each fixed depth, sufficiently long combs
have the same prefix, so $(c_k)$ is Cauchy for the classical tree ultrametric
and converges to an infinite comb in the corresponding tree completion
\cite{ArnoldNivat1980,Courcelle1983}.

By contrast, take a total extension with carrier $A\sqcup\{p,q\}$.  Preserve
all defined base applications and set
\[
 f_T(a,z)=p,\qquad f_T(p,z)=q,\qquad f_T(q,z)=p.
\]
These prescriptions are compatible because the first displayed base
application is undefined and the other two involve new points.  Fix
$a_0\in A$ and assign the value $a_0$ to every operation-table entry not
already forced either by compatibility or by the displayed prescriptions.
Then the images of $c_k$ alternate between $p$ and $q$ from $k=1$ onward.
Thus this single complement-two extension separates every pair
$c_k,c_{k+1}$ for $k\ge1$, so $(c_k)$ is not Cauchy modulo $\theta_2$.  Hence
the two uniformities are genuinely different.
\end{remark}

\section{Conclusion}

Finite-complement compatible extensions define a natural Hausdorff congruence
topology on the free compatible extension $F(D)$ of a partial algebra.  The
finite-subset separation needed for Hausdorffness is a short Rees-congruence
argument on finite subterm closures and belongs to standard finite-state
algebraic machinery.

For a finite retained base, the topology is exactly the ordinary profinite
topology.  The genuinely different situation begins with an infinite base,
where an allowed target algebra may itself be infinite although only finitely
many elements lie outside the fixed copy of $A$.

The main theorem gives a sufficient condition for incompleteness in that
setting.  Starting from an undefined base application, a nontrivial unary
context whose orbits on the base have uniformly bounded finite cardinality
produces an orbit whose image in every extension with $m$ new elements has a
uniform finite bound depending only on $m$.  Factorial sampling therefore gives
a common Cauchy subsequence across all such extensions.  A Rees quotient
adapted to any proposed term limit separates that term from the late orbit, so
the Hausdorff completion contains a point not represented by a finite normal
term.

The hypothesis is genuinely additional to partiality.  An infinite algebra
with a nowhere-defined binary operation provides the opposite behaviour: every
pair of terms can already be separated inside the retained base itself, so
$\theta_1$ is equality and the induced topology is discrete and complete.
Thus no implication of the form ``infinite and partial implies incomplete'' is
available.

The criterion is sufficient rather than necessary.  A natural next problem is
to characterize incompleteness of the finite-complement congruence topology for
arbitrary infinite partial bases, and in particular to determine what dynamical
conditions on contexts replace the uniformly finite-orbit hypothesis.

\section*{Acknowledgements}
Generative AI tools, including OpenAI ChatGPT and Anthropic Claude, were used
during the development of this work for exploratory mathematical discussion,
proof stress-testing, literature-search assistance, assistance with Lean~4
formalization, and manuscript editing.  The author reviewed the mathematical
statements, proofs, references, and final manuscript and accepts full
responsibility for the content.

\section*{Declarations}

\paragraph{Ethical approval.}
Not applicable.

\paragraph{Competing interests.}
The author has no relevant financial or non-financial interests to disclose.

\paragraph{Authors' contributions.}
Adrian Puha is the sole author and accepts responsibility for the conception,
mathematical claims, interpretation, literature review, supplementary
formalization, and preparation of the manuscript.

\paragraph{Availability of data and materials.}
No datasets were generated or analysed for this study.  The Lean~4
formalization and supporting source files are publicly available at
\url{https://github.com/Puhyuhy/ResolutionSemantics}.

\paragraph{Funding.}
No funds, grants, or other support was received.

\begin{thebibliography}{99}

\bibitem{ArnoldNivat1980}
Arnold, A., Nivat, M.: The metric space of infinite trees: algebraic and
topological properties. Fundam. Inform. 3(4), 445--475 (1980).
\url{https://doi.org/10.3233/FI-1980-3405}

\bibitem{BergstraHirshfeldTucker2009}
Bergstra, J.A., Hirshfeld, Y., Tucker, J.V.: Meadows and the equational
specification of division. Theor. Comput. Sci. 410(12--13), 1261--1271 (2009).
\url{https://doi.org/10.1016/j.tcs.2008.12.015}

\bibitem{BulmanFleming1971}
Bulman-Fleming, S.: Congruence topologies on universal algebras. Math. Z. 119,
287--289 (1971). \url{https://doi.org/10.1007/BF01109881}

\bibitem{Burmeister1986}
Burmeister, P.: A Model Theoretic Oriented Approach to Partial Algebras:
Introduction to Theory and Application of Partial Algebras---Part I.
Mathematical Research, vol. 32. Akademie-Verlag, Berlin (1986).
\url{https://doi.org/10.1515/9783112720875}

\bibitem{CainMaltcev2013}
Cain, A.J., Maltcev, V.: For a few elements more: a survey of finite Rees index
(2013). \url{https://doi.org/10.48550/arXiv.1307.8259}

\bibitem{Carlstrom2004}
Carlstr\"om, J.: Wheels---on division by zero. Math. Struct. Comput. Sci. 14(1),
143--184 (2004). \url{https://doi.org/10.1017/S0960129503004110}

\bibitem{ChenAdamekMiliusUrbat2016}
Chen, L.-T., Ad\'amek, J., Milius, S., Urbat, H.: Profinite monads, profinite
equations, and Reiterman's theorem. In: Foundations of Software Science and
Computation Structures, Lecture Notes in Computer Science, vol. 9634,
pp. 531--547. Springer (2016).
\url{https://doi.org/10.1007/978-3-662-49630-5_31}

\bibitem{ClimentVidalCosmeLlopez2023}
Climent Vidal, J., Cosme Ll\'opez, E.: Functoriality of the Schmidt
construction. Log. J. IGPL 31(5), 822--893 (2023).
\url{https://doi.org/10.1093/jigpal/jzac048}

\bibitem{ComonEtAl2008}
Comon, H., Dauchet, M., Gilleron, R., L\"oding, C., Jacquemard, F., Lugiez, D.,
Tison, S., Tommasi, M.: Tree Automata Techniques and Applications. Electronic
monograph (2008). \url{https://inria.hal.science/hal-03367725}

\bibitem{Courcelle1983}
Courcelle, B.: Fundamental properties of infinite trees. Theor. Comput. Sci.
25(2), 95--169 (1983). \url{https://doi.org/10.1016/0304-3975(83)90059-2}

\bibitem{DaveyGouveiaHaviarPriestley2011}
Davey, B.A., Gouveia, M.J., Haviar, M., Priestley, H.A.: Natural extensions and
profinite completions of algebras. Algebra Univers. 66(3), 205--241 (2011).
\url{https://doi.org/10.1007/s00012-011-0155-y}

\bibitem{Diaconescu2009}
Diaconescu, R.: An encoding of partial algebras as total algebras. Inf. Process.
Lett. 109(23--24), 1245--1251 (2009).
\url{https://doi.org/10.1016/j.ipl.2009.09.008}

\bibitem{MillerOReillyQuickRuskuc2022}
Miller, C., O'Reilly, G., Quick, M., Ru\v{s}kuc, N.: On separability finiteness
conditions in semigroups. J. Aust. Math. Soc. 113(3), 402--430 (2022).
\url{https://doi.org/10.1017/S1446788721000124}

\bibitem{Nesin2011}
Nesin, G.A.R.: Completing Partial Algebra Models of Term Rewriting Systems.
Master of Logic thesis, Institute for Logic, Language and Computation,
University of Amsterdam (2011).

\bibitem{RuskucThomas1998}
Ru\v{s}kuc, N., Thomas, R.M.: Syntactic and Rees indices of subsemigroups.
J. Algebra 205(2), 435--450 (1998).
\url{https://doi.org/10.1006/jabr.1997.7392}

\bibitem{Tichy1981}
Tichy, R.F.: The Rees congruence in universal algebras. Publ. Inst. Math.
(Beograd) 29(43), 229--239 (1981).

\end{thebibliography}

\end{document}
